\section{Giới thiệu}
\label{sec:intro}

\subsection{Đặt vấn đề}

Các sàn thương mại điện tử có hàng chục nghìn đến hàng triệu sản phẩm, trong khi mỗi người dùng chỉ xem và mua một phần rất nhỏ. Hệ thống khuyến nghị giúp thu hẹp lựa chọn bằng cách đề xuất một danh sách ngắn những sản phẩm người dùng có khả năng quan tâm. Trong thực tế, dữ liệu sẵn có thường là \emph{phản hồi ngầm}: hệ thống chỉ biết người dùng đã tương tác với sản phẩm nào, chứ không có điểm số cho mọi sản phẩm.

Dữ liệu này có hai khó khăn. Thứ nhất, nó cực kỳ thưa: trong bộ dữ liệu Amazon Electronics của dự án, mỗi người dùng trung bình chỉ có khoảng 8,6 tương tác trên hơn 44 nghìn sản phẩm. Thứ hai, người dùng mới hoặc người dùng có lịch sử thay đổi liên tục; một mô hình gán vector riêng cho mỗi người dùng phải huấn luyện lại mới phục vụ được họ.

Báo cáo này xét bài toán khuyến nghị top-$K$: từ lịch sử tương tác của một người dùng, xếp hạng các sản phẩm chưa tương tác và chọn ra $K$ sản phẩm có điểm cao nhất. Trọng tâm là các mô hình biểu diễn người dùng hoàn toàn bằng lịch sử, để vừa khai thác được cấu trúc giữa các sản phẩm, vừa phục vụ được lịch sử mới mà không cần huấn luyện lại.

\subsection{Mục tiêu và phạm vi nghiên cứu}

Dự án đề xuất hai mô hình thuộc hai nhóm phương pháp khác nhau. Graph-Regularized Neighborhood Regression (GRNR) thuộc nhóm giải trực tiếp và được so sánh với EASE và GF-CF. Diffusion-Filtered History Factorization (DFHF) thuộc nhóm cần huấn luyện và được so sánh với LightGCN, XSimGCL và DirectAU. Hai mô hình được đánh giá độc lập, không so sánh với nhau. Báo cáo trả lời bốn câu hỏi:
\begin{enumerate}[label=\textbf{RQ\arabic*.}, leftmargin=*]
\item Trong nhóm của mình, mỗi mô hình đề xuất xếp hạng tốt đến đâu so với các baseline cùng nhóm?
\item Khi chỉ giữ lại 75\%, 50\% hoặc 25\% tương tác huấn luyện, chất lượng của mỗi mô hình thay đổi thế nào so với baseline cùng nhóm?
\item Chi phí tính toán (thời gian huấn luyện, suy luận, dung lượng lưu trữ) của mỗi mô hình so với nhóm của nó ra sao?
\item Hệ thống phục vụ lịch sử mới và giải thích gợi ý như thế nào?
\end{enumerate}

Phạm vi đánh giá được giới hạn rõ. Mọi kết quả đo trên tập validation của một bộ dữ liệu (Amazon Electronics), với một seed và cho những người dùng có sản phẩm mục tiêu nằm trong tập ứng viên. Chất lượng được đo bằng thước đo xếp hạng offline; báo cáo không đánh giá doanh thu, mức hài lòng hay hành vi của người dùng thật. Vì các câu hỏi được đặt sau khi đã có kết quả thực nghiệm, mọi phân tích đều mang tính khám phá.

\subsection{Đóng góp của dự án}

Dự án có ba đóng góp:
\begin{enumerate}
\item \textbf{Hai mô hình biểu diễn người dùng bằng lịch sử.} GRNR giải một hồi quy nhỏ có điều chuẩn đồ thị cho mỗi sản phẩm, không cần vòng lặp huấn luyện. DFHF học hai bảng embedding sản phẩm được lọc khuếch tán, huấn luyện bằng softmax trên toàn danh mục. Cả hai đều chấm điểm lịch sử mới mà không cần huấn luyện lại.
\item \textbf{Đánh giá có kiểm soát trong từng nhóm.} Mỗi mô hình được so với các baseline cùng nhóm trên cùng dữ liệu, cùng tập ứng viên và cùng nhóm người dùng, ở bốn mức giảm tương tác, kèm kiểm định bootstrap ghép cặp. Báo cáo trình bày cả kết quả bất lợi và đối chiếu với một mốc xếp hạng theo độ phổ biến.
\item \textbf{Hệ thống có thể truy nguyên và demo.} Mọi số liệu trong báo cáo truy được về tệp kết quả và script sinh ra chúng; ứng dụng demo cho phép thử gợi ý với lịch sử tùy ý và xem giải thích của GRNR.
\end{enumerate}

Phần còn lại của báo cáo được tổ chức như sau. Chương~\ref{sec:background} trình bày cơ sở lý thuyết và hai nhóm phương pháp liên quan. Chương~\ref{sec:method} mô tả chi tiết hai mô hình đề xuất. Chương~\ref{sec:experiments} trình bày dữ liệu, giao thức, kết quả của từng nhóm, hệ thống demo và các hạn chế. Chương~\ref{sec:conclusion} kết luận và nêu hướng phát triển.
